\chapter*{Abstract}

Container orchestration platforms such as Kubernetes and HashiCorp Nomad have become foundational infrastructure for managing distributed workloads, yet they impose substantial operational overhead and require specialized expertise to maintain. Their predominant push-based scheduling model introduces fault-tolerance challenges, including stale delivery states and potential workload duplication under partial failures. This thesis explores a pull-based coordination model in which worker nodes independently poll a lightweight control plane for available workloads and acquire exclusive leases to claim ownership, delegating coordination to a relational database rather than implementing custom scheduling algorithms, consensus protocols, or persistent agent state. A prototype is implemented and evaluated on distributed infrastructure across scenarios covering concurrent scheduling, automatic failure recovery, and network routing. The results demonstrate that the pull-based model achieves reliable workload placement and fault recovery while requiring a fraction of the resources consumed by existing platforms, suggesting that a simpler approach to orchestration can remain practical without sacrificing core reliability guarantees.